\documentclass[journal]{IEEEtran}
% IEEE Transactions on Industrial Informatics -- anonymous regular-paper draft
% Compile: latexmk -pdf main.tex
\usepackage{amsmath,amssymb,amsfonts}
\usepackage{algpseudocode}
% Local XeTeX/Tectonic builds: route IEEEtran's ptm (Times) family requests to
% the legacy Times Type1 faces so metrics match the pdfLaTeX reference build;
% the latexmk -pdf (pdfLaTeX) path is unaffected.
\usepackage{iftex}
\ifxetex
  \DeclareFontFamily{TU}{ptm}{}
  \DeclareFontShape{TU}{ptm}{m}{n}{<->ptmr8t}{}
  \DeclareFontShape{TU}{ptm}{b}{n}{<->ptmb8t}{}
  \DeclareFontShape{TU}{ptm}{m}{it}{<->ptmri8t}{}
  \DeclareFontShape{TU}{ptm}{b}{it}{<->ptmbi8t}{}
  \DeclareFontShape{TU}{ptm}{m}{sl}{<->ptmro8t}{}
  \DeclareFontShape{TU}{ptm}{m}{sc}{<->ptmr8t}{}
\fi
\usepackage{graphicx}
\usepackage{booktabs}
\usepackage{tabularx}
\usepackage{array}
\usepackage[table]{xcolor}
\usepackage{stfloats}
\usepackage{placeins}
\usepackage{url}
\urlstyle{same}
% Cross-references and citations. The cite package renders IEEE-style citation
% groups ([14]--[16], [20], [21]) with hyperref support.
\usepackage{cite}
\usepackage[colorlinks=true,linkcolor=black,citecolor=black,urlcolor=black]{hyperref}
\hypersetup{bookmarksnumbered=true,
  pdftitle={AgentWorkShop: A Modular Task-to-Device Integration Framework for Industrial Agents},
  pdfauthor={},
  pdfsubject={IEEE Transactions on Industrial Informatics},
  pdfkeywords={industrial informatics, industrial agents, multi-agent systems, production system integration, semantic data integration}}

% Float placement only; IEEEtran margins and body type remain unchanged.
\setcounter{topnumber}{4}
\setcounter{bottomnumber}{3}
\setcounter{totalnumber}{7}
\setcounter{dbltopnumber}{2}
\renewcommand{\topfraction}{0.95}
\renewcommand{\bottomfraction}{0.80}
\renewcommand{\textfraction}{0.03}
\renewcommand{\floatpagefraction}{0.72}
\renewcommand{\dbltopfraction}{0.95}
\renewcommand{\dblfloatpagefraction}{0.72}

% Tight but conventional float separation; body type and margins are unchanged.
\setlength{\textfloatsep}{8pt plus 2pt minus 3pt}
\setlength{\dbltextfloatsep}{8pt plus 2pt minus 3pt}
\setlength{\intextsep}{7pt plus 2pt minus 3pt}
\setlength{\abovecaptionskip}{3pt plus 1pt minus 1pt}
\setlength{\belowcaptionskip}{2pt plus 1pt minus 1pt}

\input{resources/results-macros}

\newcommand{\system}{the platform}

\begin{document}
\bstctlcite{IEEEbstctl}

\title{AgentWorkShop: A Modular Task-to-Device Integration Framework for Industrial Agents}

\author{}
\markboth{IEEE Transactions on Industrial Informatics, Submitted Manuscript}{AgentWorkShop: Task-to-Device Integration}

\maketitle

\begin{abstract}
Industrial agents need a framework that connects task execution to
production context, operations, and device interfaces. AgentWorkShop routes
registered-harness tasks and direct calls through a common operation boundary
to context checks, semantic I/O, configured drivers, and service-owned
records. This boundary separates tool discovery from execution checks and
keeps device receipts distinct from later process observations. We evaluated
the registered \texttt{omp} harness adapter, the Pi CLI adapter, and a direct
Representational State Transfer (REST) route in cast-film and annealing simulators with
different signal and protocol mappings. Six sequential \texttt{omp} tasks and
two Pi CLI fixture tasks completed with synthetic approval, readback-verified
receipts, and linked run-tagged samples; direct REST calls retained
operation records. All reported samples remained outside target bands.
Mock-state tests refused 12 stale-context changes before writes; a focused
ablation let six product/run changes reach mock writes. These results
demonstrate selected route integration and one check's effect, but not process
improvement, general harness interchangeability, or comparative platform
benefit.
\end{abstract}

\begin{IEEEkeywords}
Industrial agents, industrial communication, industrial informatics,
production systems, semantic data integration.
\end{IEEEkeywords}

\IEEEpeerreviewmaketitle

\section{Introduction}
\IEEEPARstart{I}{ndustrial-agent} deployments must connect each task to its
production context, process signals, and equipment interfaces. Existing work
addresses end-to-end automation control, operator-reviewed control-sequence
generation, modular factory skills, and production scheduling
\cite{xia2025control,huang2026made,koehle2026llmdriven,waseem2026pretrained}.
Other studies target PLC code generation, context-aware control, resource
capability discovery, shop-floor coordination, digital-twin fault handling,
natural-language machine access, and robotic-line control
\cite{liu2026agents4plc,wu2025instructmpc,lim2025llmenabled,zhao2026shopfloor,gill2025llmdt,hofmann2025nlc}.
Across production lines, agent roles, signal maps, protocols, objectives, and
product/run context vary. An integration architecture must accommodate these
differences around shared task and operation services while keeping each
action tied to its active production context. Recent industrial-agent
architectures add AAS-grounded OPC UA actions with signed approval,
P\&ID-grounded proposal checks, and governed sensor--MES--ERP evidence
streams \cite{graul2026aas,schall2026safe,lin2026agentic}. These works address
asset-specific execution, process validation, or closed-loop analytics. At
their software-integration boundary, this paper examines a specific
implementation requirement: distinct task adapters and direct tool callers
should invoke one shared named operation handler while retaining evidence
appropriate to each route. Task routes can link task, approval, command
receipt, and process observation; a direct call has operation records but no
task edge. This is a scoped implementation contract, not an assumption that
prior platforms lack equivalent behavior.

Consider two production scenes with different process signals and interfaces:
cast-film thickness control and annealing-zone temperature control. An agent
task in either scene needs production context, typed signal access, and a
route from a named operation to the configured device protocol. The framework
must also distinguish a command receipt from a subsequent process
observation. The shared task and operation services can therefore be
configured with scene-specific signal maps, protocols, and limits while
retaining a common execution path. The evaluation instantiates this path in
two simulator scenes.

This paper presents AgentWorkShop as a modular framework for task-to-device
integration in industrial agents. It connects collaborative tasks through a
registered harness interface to production context, manufacturing execution
system (MES) data, analysis and modeling services, data-acquisition (DAQ) and
write-control (DCW) nodes, and protocol-specific drivers. Its interfaces include
command-line, web, software development kit (SDK), Representational State
Transfer (REST), Model Context Protocol (MCP), and Agent-to-Agent
(A2A) surfaces. Task and chat routes enter the
collaboration runtime, whereas a direct REST tool-invocation route can call
the shared operation dispatcher without creating a task.
Scene blueprints define signal maps, objectives, and guards; channel profiles,
tool policies, and agent bindings scope operations. Agents and callers invoke
named operations through shared runtime checks; the framework does not choose
routes autonomously and retains proposals, command confirmation, and process
observations as distinct records. Fig.~\ref{fig:arch} presents the system
architecture, and Fig.~\ref{fig:runtime} shows the cross-module task flow.
Section~IV details the node-exploration and recipe lifecycles.

This paper tests an execution-and-evidence contract by asking whether distinct
task adapters and a direct-call surface reach the same named industrial
operation handler and retain route-specific records linking command and
observation.
In the evaluated paths, the \texttt{recipe\_trial} operation key resolves
through the shared dispatcher to \texttt{toolRecipeTrial}; Table~\ref{tab:routes}
lists this terminal handler for each route. Tests use registered
\texttt{omp} and Pi adapters and direct REST in two simulator scenes; broader
adapter interchangeability remains untested. A separate mock-state experiment
tests pre-dispatch refusal after an approved context changes; a deterministic
fixture replays OMP without external model decisions. The suite injects line,
product, recipe, and run changes during approval for \texttt{recipe\_apply},
\texttt{recipe\_trial}, and dispatching \texttt{recipe\_rollback}; all stale
actions are refused before mock-driver writes. A post-recheck fault injection
also tests a context change queued before the dispatch continuation. In a
focused test-only ablation that omits only the tuple-equality check, six
stale product/run changes reach mock actuation, while independent checks still
refuse the other six cases.

The paper makes two contributions:
\begin{enumerate}
\item It specifies and implements a modular execution-and-evidence contract
  that routes registered-harness tasks and direct calls to the same named
  operation handler through common checks, semantic I/O, configured protocol
  drivers, and
  service-owned records. Task routes retain task-to-approval-to-receipt-to-
  observation links where available; direct calls retain operation records
  without creating a task. For approved recipe actions, the
  implementation binds approval to a line--product--recipe--run tuple and
  rechecks it after asynchronous authorization validation, immediately before
  entering software dispatch. This is a software-state guard, not an atomic
  lock across device I/O. It
  separates prompt grounding, dispatch authority, device receipt, and later
  process observation; it is an implementation contribution, not a new
  protocol, safety-verification method, or process-control algorithm.
\item It evaluates route-specific record continuity for the OMP and Pi task
  adapters and direct REST calls across cast-film and annealing simulator
  mappings, with each route reported separately. It also reports a
  deterministic OMP replay and a mock-state ablation of the context check.
  These measurements test selected route compositions, not process
  attainment, general adapter portability, or comparative platform value.
\end{enumerate}

Section~II reviews related industrial-agent and integration work.
Section~III presents the design rationale, architecture, and cross-module
execution flow; Section~IV explains module responsibilities and their
integrated operation. Section~V describes the experiments and results, and
Sections~VI and~VII discuss the evidence limits and conclude the paper.

\section{Related Work}
\label{sec:related}
This section relates AgentWorkShop to plant-facing industrial agents and to
production and automation abstractions for context, tool access, and
operational constraints. Table~\ref{tab:nearest} compares integration and
execution features described in selected industrial-agent architectures,
including systems with adjacent but distinct process-control objectives. The
comparison is focused rather than exhaustive; a feature marked as not
specified was not described in the cited study and is not inferred to be
absent.
Studies were included in this comparison when they describe an industrial
agent's production context, plant-facing operations, or an action, approval,
or evidence pathway. This criterion defines the table's scope; it is not a
comprehensive survey or a performance ranking.
Feature entries were coded from the cited study's full text only when the
feature was explicitly described; otherwise the cell is marked ``Not
specified.'' Evaluation descriptions are transcribed at the level reported
by each source, without normalizing sample units, success criteria, or
metrics across studies.

\subsection{Industrial Agents and Plant Integration}
Industrial-agent systems, including large language model (LLM)-based systems,
span end-to-end automation control, supervisory
services, programmable logic controller (PLC) interfaces, manufacturing
coordination, and operator-facing machine access. Xia et al. present an
LLM agent framework for end-to-end industrial automation control, using
structured prompts and event-driven context for production planning and
operation \cite{xia2025control}. Huang et al. present a multimodal
architecture whose stated focus is automated control generation for
industrial cyber-physical systems \cite{huang2026made}. Biskupovic et al.
study agent-based supervision for service-oriented industrial systems
\cite{biskupovic2025supervision};
Agents4PLC and InstructMPC address PLC code generation and human--LLM
context-aware control, respectively
\cite{liu2026agents4plc,wu2025instructmpc}. Manufacturing studies consider
dynamic resource-capability exploration \cite{lim2025llmenabled}, multi-agent
shop-floor coordination \cite{zhao2026shopfloor}, and digital-twin-supported
fault handling \cite{gill2025llmdt}. Natural-language machine access
\cite{hofmann2025nlc} and real-time LLM control of a robot-operated serial
line \cite{waseem2026pretrained} extend this work to operator interfaces and
robotic-line control. Together, these studies address different operational
contexts, interfaces, and allocations of responsibility between agents,
operators, and controllers.

The closest architecture-level comparison is Xia et al., who connect cognitive,
digital-twin, and physical-automation layers and organize twin information as
state snapshots, event messages, and plan sequences; their case studies
address robotic automation and process simulation
\cite{xia2025architecture}. Other nearby designs expose factory-module
operations through MCP with state-based tool filtering, plan constrained
manufacturing steps, or compose semantic recipes from OPC UA skills
\cite{koehle2026llmdriven,silva2025mcpmanufacturing,koehle2025semanticrecipes}.
These studies describe important parts of plant-facing integration. The
recent systems in Table~\ref{tab:nearest} also combine agent execution with
approval, industrial data, or audit records. Graul et al. ground proposals in
the Asset Administration Shell (AAS) and route approved actions through OPC UA
Methods, with signed approvals and audit logging
\cite{graul2026aas}. Schall separates advisory language-model agents from
process and instrumentation diagram (P\&ID)-grounded verification and
execution, and evaluates structural checks across five simulated process
scenarios \cite{schall2026safe}; its context includes supervisory control and
data acquisition (SCADA) and standard operating procedure (SOP) information.
Lin and Chen build a governed sensor--MES--ERP digital thread for closed-loop
monitoring and control, evaluated on real manufacturing
records \cite{lin2026agentic}. These works address asset-specific approval,
process validation, or closed-loop analytics. AgentWorkShop examines the
complementary software-integration question of whether distinct registered
task adapters and a direct-call surface can reach one named industrial
operation through a common handler while retaining route-appropriate records.
The evaluated contract covers production context, semantic I/O, configured
drivers, command receipts, and later process observations. Its route-specific
approval check also rechecks the line--product--recipe--run tuple immediately
before software dispatch. This is a scoped implementation contract, not a
claim that other systems lack equivalent mechanisms. The study tests the
contract with two registered adapters, direct REST, and two simulator
mappings; it does not establish general adapter interchangeability, reduced
integration effort, or superiority over another platform.
Table~\ref{tab:nearest} summarizes only
features explicitly described in the cited studies.
\begin{table*}[!t]
\centering
\caption{Feature-level comparison of selected plant-facing integration
studies. Entries summarize reported design scope and evaluation, not performance.}
\label{tab:nearest}
\small
\setlength{\tabcolsep}{3pt}
\begin{tabularx}{\textwidth}{@{}>{\raggedright\arraybackslash}p{1.45cm}>{\raggedright\arraybackslash}p{3.85cm}>{\raggedright\arraybackslash}p{2.7cm}>{\raggedright\arraybackslash}X@{}}
\toprule
System & Context and operation checks & Device-facing route & Reported evaluation \\
\midrule
Xia et al. \cite{xia2025architecture} & Cognitive, twin, and physical layers; snapshots, events, plans & Digital-twin/physical automation & Robotic automation and process-simulation cases \\
K{\"o}hle et al. \cite{koehle2026llmdriven} & Factory agents with injected live state; state-based tool filtering & MCP exposes OPC UA methods; MQTT coordination & Nine simulated production challenges \\
Vieira da Silva et al. \cite{silva2025mcpmanufacturing} & Constrained LLM planning; operation checks not specified & MCP-exposed manufacturing functions & Laboratory multi-step process \\
K{\"o}hle et al. \cite{koehle2025semanticrecipes} & AI-assisted recipes from OPC UA skills; checks not specified & OPC UA skills and virtual commissioning & Unity assembly-workstation proof of concept \\
Graul et al. \cite{graul2026aas} & AAS-grounded context; signed human approval & PLC-validated OPC UA Methods; immutable audit log & Two laboratory PoCs; single scripted runs \\
Schall \cite{schall2026safe} & P\&ID, SCADA, time-series, SOP; topology and rule checks & Advisory proposal with gated execution & Five simulated scenarios; 43 proposals and 50 LLM runs \\
Lin and Chen \cite{lin2026agentic} & Sensor--MES--ERP thread; SOP-, risk-, capability-aware gating & Edge monitoring and closed-loop actions & 25,275 real records; operational metrics \\
\midrule
This work & OMP and Pi registered task adapters plus direct calls; approval-bound context rechecked before recipe dispatch & Shared dispatcher; semantic nodes resolve to configured DAQ/DCW drivers & Three tested entry paths across two simulator mappings; deterministic Pi fixture; no attainment claim \\
\bottomrule
\end{tabularx}
\end{table*}

The operation checks complement safety-aware decision support
\cite{jiang2026agentic} but do not replace independent process protection.

\subsection{Production Context, Agent Interfaces, and Automation}
Industrial integration also depends on abstractions that are not specific to
agent architectures. The Model Context Protocol (MCP) specifies interoperable
tool discovery and call mechanisms \cite{mcp2024}; the operation-level scope,
policy, grant, and binding checks discussed here are AgentWorkShop's own
dispatcher rules, distinct from MCP transport authorization. In parallel,
ISA-88 and ISA-95 describe recipes and production operations
\cite{isa88_2010,isa95_2025}. Run-to-run methods connect process observations
to successive production runs \cite{moyne2000r2r}. At the agent level,
AgentSpec and CaMeL address tool behavior and capability-based information
flow \cite{wang2025agentspec,debenedetti2025camel}. Together, these works
provide complementary abstractions for production context, tool interaction,
and information flow. AgentWorkShop composes such abstractions with its task
runtime, operation checks, semantic I/O, and route-specific records; this
paper evaluates selected paths through that composition rather than proposing
a replacement for the underlying standards.

\begin{figure*}[!t]
\centering
\includegraphics[width=\textwidth]{resources/fig-architecture-2026-9pt.pdf}
\caption{AgentWorkShop architecture and record ownership. Task/chat calls pass
through the listed access surfaces and a registered harness; direct REST and
selected stdio-MCP operation calls enter the shared dispatcher, while MCP
management/task methods use their service routes. R22 evaluates direct REST
operation calls; the stdio-MCP route is shown but not evaluated. Recipe writes recheck the
approved line--product--recipe--run tuple; each service retains its own
records. Independent programmable logic controller and safety-instrumented
system protection remains external.}
\label{fig:arch}
\end{figure*}
\section{System Design}
\label{sec:design}
AgentWorkShop is designed for industrial tasks whose agents, production
contexts, interfaces, and process signals vary between deployments. Its
architecture organizes these variations around shared task and operation
services, checks each proposed operation against the active production
context before device routing, and retains route-specific evidence that
distinguishes a command receipt from a measured process response.

\subsection{Design Rationale and Requirements}
\label{sec:principles}
Five requirements---scene variability, harness variability, operation
authority, outcome observability, and extensibility---organize six design
rules (D1--D6). These are deployment constraints drawn from the target
workflows. D1 stores scene settings as versioned data;
D2 provides registered harness, tool, and driver interfaces; D3 separates
proposal from execution; D4 treats the tool catalog as discovery rather than
permission; D5 routes actuator writes through the common DCW admission point;
and D6 retains evidence by operation.

The requirements trace to implementation as follows.
Scene variability maps to D1, scene blueprints, and context resolution;
harness variability and extensibility map to D2, the registered harness and
driver interfaces; operation authority maps to D3--D5, the shared dispatcher,
grants, bindings, and write admission; outcome observability maps to D6,
route-specific receipts, and observations. These interfaces are design
extension points; the evaluation in Section~\ref{sec:eval} tests selected
registered-harness and direct-operation routes, not general harness
interchangeability. AgentWorkShop supplies task and operation services;
agent strategy and independent process protection remain outside this design.

\subsection{Layered Architecture}
\label{sec:layers}
\label{sec:arch}
Fig.~\ref{fig:arch} summarizes the access routes, shared dispatch path,
production services, semantic I/O, configured drivers, and service ownership
of records. Task and chat routes enter the
collaboration band, which contains workspaces, channels, task queues, shared
memory, mailboxes, task templates, and harness adapters. Direct REST tool
calls and the stdio MCP operation bridge can call the shared dispatcher
without creating a task or invoking a harness; MCP management and task methods
remain on their service routes. Direct calls have no agent-turn prompt
context, so handlers resolve current route state from the request and services
at invocation. The governance band projects an operation catalog and
checks each invocation. The context band assembles bound-line and active-run
information for prompt grounding, while operation handlers resolve the live
state relevant to their route. The service band provides acquisition, write
control, analysis, retrieval, and optional model services. The field-I/O band
resolves semantic nodes to built-in or plugin drivers for a plant or
simulator.

Each band has a distinct responsibility. Task/chat routes enter the task
runtime; direct operation routes enter the dispatcher; and both use shared
operation checks where applicable. A harness adapter supplies
agent-specific interaction on task paths, semantic nodes hide protocol
details from operations, and the driver registry provides the device-facing
extension point. Owning services retain task, command, observation, and
version records separately and link them to available actor and production
context. No single status conflates these outcomes.

\subsection{Cross-Module Execution Flow}
\label{sec:executionflow}
Fig.~\ref{fig:runtime} summarizes a task turn from entry through operation
execution, with route-specific evidence retained by owning services. Registered
harness adapters receive a scoped tool catalog during setup and supported
registry refreshes; the runtime assembles prompt-time production context for
each agent turn. This context is grounding information, not a task-pinned
transaction. For each returned tool call, the shared dispatcher rechecks
invocation scope and channel policy and, for industrial operations, current
grant and binding conditions before handler selection. The selected handler
then applies route-specific guards, active-run checks, approval
revalidation, and device-I/O steps as applicable. Command receipts and later
process observations remain distinct records.

For an industrial operation $u$ by agent $a$ in channel $c$, define the
shared dispatch condition as
\begin{equation}
\mathcal{A}(a,c,u)=\bigwedge_{X\in\mathcal{C}(u)}X(a,c,u),
\label{eq:dispatch-contract}
\end{equation}
where $S$, $P$, $G$, and $B$ denote invocation scope, current channel policy,
runtime grant, and the agent's binding to referenced resources, respectively.
$\mathcal{C}(u)\subseteq\{S,P,G,B\}$ is the set of predicates declared for
operation $u$; it always contains $S$ and $P$. The grant predicate $G$ is
omitted only for an explicitly grant-exempt industrial tool, and $B$ is
included when the operation declares resource references. A declared
no-reference operation omits $B$. Each included check must pass before
handler selection. The policy predicate includes channel tool selection,
profile, and duty-separation checks where configured. This condition is
necessary for handler selection, not sufficient for actuation: the selected
route applies its own guards and approval, then rechecks authority, context,
and applicable limits before dispatch.

Fig.~\ref{fig:execution} gives reference pseudocode for the registered-harness
task path, not one literal orchestrator function. Direct REST tool invocation
bypasses the task runtime and harness but enters the same shared dispatcher.
Route-specific timing, approval, and evidence rules are detailed in
Section~\ref{sec:modules}.

\begin{figure*}[!t]
\centering
\includegraphics[width=\textwidth]{resources/fig-runtime-2026-9pt.pdf}
\caption{Registered-harness task flow. Context is assembled per turn; the
scoped tool catalog is loaded at setup or refresh. The dispatcher checks
scope, policy, and applicable grants and bindings, then returns each tool
result to the active harness turn. Direct REST and selected stdio-MCP
operations enter the shared dispatcher without a task. In R20, the follow-up
probe, ingest, and run-tagged query were initiated by the test runner outside
the agent task. A command receipt remains distinct from a process observation.}
\label{fig:runtime}
\end{figure*}

\begin{figure}[!t]
\centering
\begin{algorithmic}[1]
\Require task $q$, channel $c$, bound agent $a$, registered scoped catalog $T$
\Ensure task result; any retained evidence remains with its owning service
\State $x_t \gets \Call{ResolvePromptContext}{c,q}$
\State $m \gets \Call{RetrieveMemory}{q}$
\State $p \gets \Call{AssemblePrompt}{q,x_t,m}$
\State $h \gets \Call{InvokeRegisteredHarness}{a,p,T}$
\ForAll{tool request $u$ returned by $h$}
  \State $g \gets \Call{CheckScopeAndChannelPolicy}{a,c,u}$
  \If{\textbf{not} $g$}
    \State \Call{EmitToolErrorToHarness}{h,u}; \textbf{continue}
  \EndIf
  \If{$u$ is an industrial operation}
    \State $v \gets \Call{CheckGrantBindingOwnership}{a,u}$
    \If{\textbf{not} $v$}
      \State \Call{EmitToolErrorToHarness}{h,u}
      \State \textbf{continue}
    \EndIf
  \EndIf
  \State $k \gets \Call{SelectHandler}{u}$
  \State $r \gets \Call{ExecuteRouteHandler}{k,a,c,u}$
  \Comment{handler resolves live state and may persist route-specific records}
  \State \Call{EmitDeclaredToolResultToHarness}{h,r}
\EndFor
\State \Return task state and tool results
\end{algorithmic}
\caption{Reference pseudocode for the registered-harness task path.
Direct operation invocation is not shown. Evidence persistence is
route-specific and remains with the owning service; the shared tool result
does not carry a generic evidence envelope.}
\label{fig:execution}
\end{figure}

\section{Module Functions and Integration}
\label{sec:modules}
The architecture in Fig.~\ref{fig:arch} is implemented as cooperating
modules rather than a single monolithic agent. Shared interfaces connect
task execution to configuration, industrial services, protocol drivers, and
route-specific evidence. This section describes those responsibilities and
the checks that preserve their boundaries.

The implementation scope is broader than the evaluated scope. Section~V
separates implemented interfaces from the routes and checks exercised here;
an interface intended for extension does not by itself establish
interchangeability or portability.

The extension interfaces specify the information passed between modules: a
harness call carries a tool name and arguments and returns response text with
an error flag. A DCW driver accepts an engineering value, tolerance, range,
and driver configuration; its result separates write status, raw value,
readback, and acknowledgment grade (readback-verified, transport-ack, or
unverified). Each DAQ sample records its node, timestamp, value, and state,
plus optional line, product, recipe, and run IDs. A write anchor records
line/node, prior and requested values, source, actor, time, and optional
run/task/approval IDs. Recipe definitions retain separate versions. Together,
these fields keep tool errors, device receipts, observations, and production
lineage inspectable as separate records.

DAQ and write-anchor records may omit line, product, recipe, run, task, or
approval identifiers when a route does not supply them. For the R20 trace
audit, the join keys are required: task ID, approval ID, recipe-trial run ID,
active-run ID, and the run ID returned with the sample must agree with the
corresponding route records. R21 checks the corresponding IDs on its
deterministic fixture path. Thus the reported complete joins are requirements
of those evaluated records, not a guarantee that every general record
contains every identifier.

\subsection{Task Runtime and Harness Adapters}
Task and chat access surfaces create or join a workspace channel, where the
task engine manages members, task queues, shared memory, and execution
templates. Direct REST tool calls and selected stdio MCP operation calls can
enter the shared dispatcher without creating a task or invoking a harness;
MCP management and task methods use their corresponding services. A direct
call has no agent-turn prompt context: its operation handler resolves live
route state from request identifiers and service data at invocation.
The harness registry instantiates an adapter behind a common agent interface;
the prompt builder combines the task with the active scene brief, industrial
context, channel profile, agent binding, and retrieved memory. Tool calls
return to the shared dispatcher regardless of which registered harness
produced them. The registry keeps engine-specific execution behind the common
agent interface. Industrial context is assembled for
each agent turn rather than pinned for the task lifetime, so prompt context is
grounding information rather than write authorization. An approved recipe
action separately captures the current line--product--recipe--run tuple when
approval is requested and compares it again before dispatch. This
approval-time check does not imply that the earlier prompt context remained
unchanged.

\subsection{Production Context and Configuration}
Scene blueprints describe devices, signal maps, objectives, and process
guards. A task resolves these definitions against its active
line--product--recipe--run context. The channel profile selects the available
capability family, the tool policy narrows the operation catalog, and the
agent binding scopes the nodes, controls, and recipe operations available to
that member. These filters can restrict one another but cannot broaden the
underlying device range. The dispatcher re-reads policy at execution, so a
cached catalog does not retain authority after a grant or policy change. The
prompt builder assembles the agent's available bound-line and active-run
context on each turn; this prompt-time context is grounding information and
does not pin the complete tuple to the task.

An approved recipe action separately captures the active tuple when approval
is requested. After the approval wait, the handler refuses software dispatch
if the line, product, recipe, or run field has changed. This route-level check
is separate from prompt-context assembly and does not lock physical I/O.

Specific production snapshot, trial, and MPC routes, together with Hybrid Twin
training submissions, resolve and validate a versioned \texttt{SceneContract}
at operation time. SceneContract enforcement is route-specific: the Twin
physics-specification draft route can instead accept a caller-supplied scene
description, and therefore does not validate or pin all prompt context.

\subsection{Acquisition, Analysis, and Modeling Services}
Semantic DAQ nodes expose typed values, engineering units, ranges, and
sampling settings while hiding protocol-specific reads. A per-node DAQ
runtime schedules acquisition and stores timestamped observations; the
gateway aggregates node pipelines. MES-facing operations provide the
configured production-data path. Analysis and retrieval services consume
these records, while optional model, training, optimization, and digital-twin
services are invoked through their registered operation handlers. These
services complement the agent workflow; AgentWorkShop does not prescribe a
single optimization strategy.

\subsection{Semantic I/O and Actuation}
The driver registry resolves a semantic node to a built-in or plugin driver.
DAQ handlers use the node's acquisition interface; DCW handlers route a
write or readback through the controller and the configured protocol driver.
This separation lets recipes and agents refer to semantic controls without
embedding Modbus, OPC UA, MQTT, HTTP, or MES transport details in their
operation arguments. The supported driver mappings are implementation
interfaces; only the mappings and workflows enumerated in Section~\ref{sec:eval}
are experimentally exercised.

\subsection{Shared Dispatch and Write Admission}
\label{sec:governed}
Every named operation passes the dispatcher. Its ordered checks validate
invocation scope and current channel policy (tool selection, profile, and
applicable duty separation), then the agent's binding and runtime grant and
ownership of referenced resources.
A refusal returns the excluding condition. Denial persistence is selective:
grant denials and selected resource-reference denials have explicit audit
paths, whereas scope and channel-policy denials are not uniformly retained as
durable audit events. A catalog entry therefore describes availability but
does not authorize execution. Handlers for operations that require manual
approval suspend for a decision, then revalidate authority and active limits
before dispatch. For manual recipe operations, the handler checks the active
run, candidate parameters, and route timing before approval; after the wait it
rechecks current authority and active limits. For \texttt{recipe\_apply}, a
changed definition requires a new proposal. These route-specific checks
complement, but do not replace, the per-call dispatcher checks.

All actuator writes converge on the DCW admission point. For an online node
write, the admissible interval is
\begin{equation}
\mathcal{I}^{\mathrm{online}}_n(t)=\mathcal{H}_n\cap\mathcal{B}_n
\cap\mathcal{P}_n(t)\cap\mathcal{R}_n(t),
\label{eq:admission}
\end{equation}
where $\mathcal{H}_n$ is the hard engineering range and $\mathcal{B}_n$,
$\mathcal{P}_n$, and $\mathcal{R}_n$ are the baseline, active-product, and
active-recipe limits. Missing optional limits impose no restriction, whereas
an empty intersection rejects the request and identifies the excluding
layer. Recipe dispatch applies the hard, baseline, and product limits to the
replacement recipe rather than intersecting with the previous recipe window.
The actuator path also enforces step, interval, and hold-window checks.
These admission rules constrain commands; process objectives and coupled
variable guards are assessed from later observations. The framework does
not replace independent PLC or safety-instrumented protection.

\subsection{Integrated Operation and Route-Specific Evidence}
\label{sec:tasklifecycle}
A recipe trial illustrates how the modules cooperate. The task runtime
delivers a named \texttt{recipe\_trial} call from the agent harness; the
dispatcher checks the agent's recipe binding and execution scope; the recipe
handler checks the active run, candidate parameters, and route timing; and a
manual binding requests approval before the DCW controller sends the batch
through its driver. The handler returns a per-parameter receipt. The caller
then obtains DAQ observations and can request a separate update, apply, or
rollback operation. A trial receipt confirms command handling; it does not
by itself establish process attainment or persist a new recipe version.

The routes retain different records because they answer different
questions. A driver acknowledgment reports command evidence; DAQ reports a
process observation; an exploration summary records a bounded probe; and a
recipe version records a saved definition. Their relation through actor and
production context allows an execution to be inspected without treating
these outcomes as interchangeable.

\input{sections/mechanisms}

\section{Experimental Evaluation}
\label{sec:eval}
The evaluation asks whether the shared operation services can be reached from
multiple task and direct-call routes across two simulator scenes with distinct
signal and protocol mappings, while retaining route-appropriate records. R20
is the primary result: six sequential tasks through the registered OMP
adapter, three per scene on one simulator instance per scene, with
task-to-approval-to-receipt-to-sample links. R21 is a deterministic fixture
replay of that OMP route. R22 exercises direct REST calls without a task or
harness. R23 exercises a second registered task adapter, Pi, using a
deterministic local model fixture. These evidence units are reported
separately. All use repository simulators; Supplementary File~S1 contains
their records, scripts, and trace checks.
\subsection{Evidence Units and Experimental Procedure}
\label{sec:testbed}
\label{sec:procedure}
R20 uses the registered harness adapter named \texttt{omp} in cast-film and
annealing scenes
(\texttt{omp} is the project adapter key, not a network protocol),
with three sequential task executions per scene on one simulator instance per
scene. Each execution creates a new logical line, product, recipe, active run,
channel, agent, and task on that shared scene simulator. Line start reapplies the recipe baseline
(120~rpm or 730~$^{\circ}$C), so the control read returns to baseline while
the simulated process state carries forward. These are sequential software
runs, not independent process replications. R21--R23 use fresh isolated
application/simulator environments and are reported separately from R20.

Table~\ref{tab:routes} summarizes the evaluated paths and their inference
limits; Table~\ref{tab:q8} gives the six R20 process snapshots.

For task paths, full trace conformance requires a completed task, approval, a
readback-verified receipt, and an active-run sample linked through task and
approval records. R20 and R23 meet this task-state criterion; R21 records the
task, approval, receipt, and sample links but does not assert task completion.
Direct REST is scored on operation records and has no task edge; route counts
remain separate.

\FloatBarrier
\begin{table*}[!t]
\centering
\caption{Evaluation paths, trace-contract outcomes, and inference limits.}
\label{tab:routes}
\small
\setlength{\tabcolsep}{4pt}
\begin{tabularx}{\textwidth}{@{}>{\raggedright\arraybackslash}p{1.25cm}>{\raggedright\arraybackslash}p{3.4cm}>{\raggedright\arraybackslash}X>{\raggedright\arraybackslash}p{4.0cm}@{}}
\toprule
Path & Entry route $\rightarrow$ handler & Recorded outcome & Inference limit \\
\midrule
R20 & Six OMP tasks $\rightarrow$ \texttt{toolRecipeTrial}; three per scene. & 6/6 task-to-sample joins and verified receipts; all probes outside target. & One simulator and one probe per scene/task; no attainment claim. \\
R21 & OMP fixture tasks $\rightarrow$ \texttt{toolRecipeTrial}; one per scene. & 2/2 linked task-context/approval/receipt/sample records; 31 runner checks passed. & Task completion not asserted; no external model decision or process replication. \\
R22 & Direct REST calls $\rightarrow$ \texttt{toolRecipeTrial}; one per scene, without a task. & 2/2 verified receipt/readbacks; run-tagged samples were 46.37~$\mu$m (48--52~$\mu$m target) and 675.20~$^{\circ}$C (700--760~$^{\circ}$C target), both below band; 19 checks passed. & No task-to-sample join or cross-harness claim. \\
R23 & Pi CLI tasks $\rightarrow$ \texttt{toolRecipeTrial}; one per scene. & 2/2 task-to-approval/receipt/sample links; both readbacks verified; 37 application and 28 isolation checks passed. & One probe per scene; no external LLM, process replication, or target attainment. \\
Context tests & Three recipe routes; full handler regression and separate test-only ablation. & Full regression refused 12/12 injected changes before writes. Ablation allowed six product/run changes to reach mock writes and create execution records; six line/recipe changes remained blocked. & Mock/temporary state; not a platform baseline or physical safety test. \\
\bottomrule
\end{tabularx}
\end{table*}

\subsection{Registered-Harness Process Path (R20)}
\label{sec:harness-e2e}
R20 ran the local v0.8.0 working-tree application with the registered
\texttt{omp} adapter and GLM-5.3-Flash. Supplementary File~S1 lists the source
reconstruction requirements; credentials are omitted and external model
decisions are not deterministic. Cast-film used OPC UA control and MQTT
thickness sensing; annealing used Modbus TCP for control and sensing. In each
scene the model selected one value from a prevalidated candidate set for a
bounded \texttt{recipe\_trial}, and a synthetic administrator approved it.
Each trial produced one readback-verified receipt. Twenty seconds later, the
runner probed and ingested one sample, then queried it by active-run ID through
the isolated SQLite fallback because the isolated MQTT queue and Timescale
endpoints were unreachable. All six tasks passed 76 runner assertions. The
S1 per-execution trace matrix links task, approval, trial receipt, and
active-run sample IDs; trial-run and sample run IDs differ but join through
the task and approval. These experiment-level joins do not imply that every
service record persists every link field. Task state was checked after
sampling, so completion-to-sample latency is unavailable. Table~\ref{tab:q8}
gives the values. Each pre-task value is a fresh control/process read during
logical-run initialization, separate from the post-receipt probe. Each
subsequent pre-task read followed the preceding post-trial probe by about
1.0--1.1~s; the simulator state persisted while each new line reapplied its
baseline. These are successive snapshots, not replicate pairs. All six
samples were outside target; one probe per task does not measure process
response or normal online DAQ.

\begin{table}[!t]
\centering
\caption{Bounded recipe-trial settings and post-readback simulator probes. All
six samples remained outside their target bands.}
\label{tab:q8}
\footnotesize
\setlength{\tabcolsep}{2.5pt}
\begin{tabularx}{\columnwidth}{@{}lXcX@{}}
\toprule
Case & Pre-task value / target band & Control: start$\to$trial & Sample \\
\midrule
Film-1 & $55.82/48$--$52~\mu$m & $120\to125$~rpm & $45.69~\mu$m \\
Film-2 & $45.69/48$--$52~\mu$m & $120\to130$~rpm & $46.69~\mu$m \\
Film-3 & $46.69/48$--$52~\mu$m & $120\to125$~rpm & $45.69~\mu$m \\
Anneal-1 & $675.40/700$--$760~^{\circ}$C & $730\to745$~$^{\circ}$C & $679.60~^{\circ}$C \\
Anneal-2 & $678.30/700$--$760~^{\circ}$C & $730\to745$~$^{\circ}$C & $677.70~^{\circ}$C \\
Anneal-3 & $676.90/700$--$760~^{\circ}$C & $730\to745$~$^{\circ}$C & $678.10~^{\circ}$C \\
\bottomrule
\end{tabularx}
\end{table}

\subsection{Deterministic Fixture Replay (R21)}
R21 used a deterministic transport fixture for the registered \texttt{omp}
adapter, without an LLM decision. One task per scene exercised approval,
verified receipt/readback, and a run-tagged sample; all 31 runner checks,
including cleanup, passed. The fixture did not assert task completion. This is
a replay of the same harness route, not an independent agent or process
replication, and is reported separately from R20.

\subsection{Direct Operation Endpoint (R22)}
R22 used a fresh isolated environment. A worker identity issued one direct
\texttt{recipe\_}\allowbreak\texttt{trial} call per scene through
\url{/api/workshop/agent-tools/invoke}, without a task, harness transport,
or external model. Synthetic approval produced a 1/1 verified receipt and
matching readback for each call. After 20~s, the runner queried one
active-run-tagged sample per scene; Table~\ref{tab:routes} reports both
values and target bands. All 19 checks passed.

\subsection{Second Registered-Harness Path (R23)}
R23 exercised the registered Pi CLI adapter with a deterministic loopback
OpenAI-compatible fixture that scripted one \texttt{recipe\_trial} call and
then \texttt{complete\_task} per scene; no external model provider was called.
Synthetic approval yielded a readback-verified receipt, completed task,
and active-run-tagged sample in each scene. After 20~s, cast-film moved from
120 to 125~rpm and sampled 46.38~$\mu$m (48--52~$\mu$m target); annealing
moved from 730 to 735~$^{\circ}$C and sampled 674.40~$^{\circ}$C
(700--760~$^{\circ}$C target). Both samples were below target. All 37
application and 28 isolation checks passed, including task-to-approval, trial,
active-run, and sample links. This tests Pi route integration, not Pi reasoning,
independent process response, or general adapter interchangeability.

\subsection{Software Contract Checks}
The local unit suite passed 566/566 tests across seven suites. In the 20-case
approval suite, twelve injected line, product, recipe, or run changes across
three recipe routes were refused before mock writes or new execution records.
One queued post-recheck context change was also refused before dispatch. In a
test-only ablation that removed tuple equality but retained authority and
limit checks, six product/run changes caused twelve mock writes and six
execution records; independent checks still blocked line/recipe changes.
This component ablation isolates the tuple check's incremental software-state
coverage; it is not a platform comparison or safety evaluation. Tests use mock or
temporary state and do not establish atomicity during device I/O or test
end-to-end DAQ tag failures.

\section{Discussion}
\label{sec:limits}
OMP, Pi-fixture, and REST routes completed the evaluated integration traces
in two simulator scenes. This case study complements asset-specific approval
and process-validation studies \cite{graul2026aas,schall2026safe}, but
measures neither process performance nor comparative platform value.

One simulator and one probe per task preclude independent process
replication. Only reported adapter configurations were tested; general
interchangeability, model behavior, online DAQ, configuration reuse, physical
transfer, and a platform baseline remain untested. The tuple check guards
software dispatch, not device I/O. Synthetic approval does not assess operator
performance \cite{parasuramanmanzey2010complacency,parasuraman2000model,bainbridge1983ironies};
operation restrictions do not establish security
\cite{greshake2023injection,liu2024pinjection,iec62443-3-3} or replace PLC/SIS
safeguards \cite{iec61511}. Future work needs a defined comparison and
independent simulator/plant trials.

Scope and channel-policy refusals are not uniformly persisted; auditability
is limited to the records and denial classes described in Section~IV.

\section{Conclusion}
\label{sec:conclusion}
AgentWorkShop unifies registered-task and direct-call access to shared
industrial services while retaining route-specific records. Tests across two
simulator mappings verify selected traces and one pre-dispatch context guard,
not process improvement or comparative platform benefit.


% Identity-neutral acknowledgment: documents AI assistance without exposing authors.
\section*{Acknowledgment}
OpenAI ChatGPT revised Sections~I--VI and drafted Python for Figs.~1, 2, and
4. OpenAI image generation supplied concepts for Figs.~1--2; final vectors
were redrawn and checked against the implementation. No AI system generated
experimental data; the authors verified the manuscript.

\bibliographystyle{IEEEtran}
% Keep IEEEtran's default bibliography size and leading.
\bibliography{resources/refs}

\end{document}
